\par\Needspace{8\baselineskip}\bigskip
\pdfbookmark[3]{Exposition, 10}{pallas-II-article10}
\begin{center}10.\end{center}
\noindent\hspace*{\parindent}Différentiant les expressions de \(A,B,C\) (4, 5, 6), on trouve
\[
\begin{aligned}
\mathrm{d}A&=(z\eta-y\zeta)\mathrm{d}t\\
\mathrm{d}B&=(x\zeta-z\xi)\mathrm{d}t\\
\mathrm{d}C&=(y\xi-x\eta)\mathrm{d}t.
\end{aligned}
\]
\par On tire des équations 45—47
\[
\begin{aligned}
-B\mathrm{d}A+A\mathrm{d}B&=mp\sin i^2\mathbin{.}\mathrm{d}\AscendingNode\\
AC\mathrm{d}A+BC\mathrm{d}B-(AA+BB)\mathrm{d}C&=(mp)^{\frac32}\sin i\mathbin{.}\mathrm{d}i\\
A\mathrm{d}A+B\mathrm{d}B+C\mathrm{d}C&=\tfrac12m\mathrm{d}p.
\end{aligned}
\]
\par Substituant ici pour \(\mathrm{d}A,\mathrm{d}B,\mathrm{d}C\) les valeurs qu’on vient de trouver, et faisant attention que \(Ax+By+Cz=0\), on obtient
\[
\begin{aligned}
-z(A\xi+B\eta+C\zeta)\mathrm{d}t&=mp\sin i^2\mathbin{.}\mathrm{d}\AscendingNode\\
(Bx-Ay)(A\xi+B\eta+C\zeta)\mathrm{d}t&=(mp)^{\frac32}\sin i\mathbin{.}\mathrm{d}i\\
\{(Cy-Bz)\xi+(Az-Cx)\eta+(Bx-Ay)\zeta\}\mathrm{d}t&=\tfrac12m\mathrm{d}p.
\end{aligned}
\]
\par Donc, substituant pour \(z\) et \(Bx-Ay\) leurs valeurs (éq. 44, 50), et introduisant les forces \(W,V\), on a
\[
\begin{aligned}
\BareRoot m\mathbin{.}r\sin(v-\AscendingNode)\mathbin{.}W\mathrm{d}t&=\BareRoot p\mathbin{.}\sin i\mathbin{.}\mathrm{d}\AscendingNode\\
\BareRoot m\mathbin{.}r\cos(v-\AscendingNode)\mathbin{.}W\mathrm{d}t&=\BareRoot p\mathbin{.}\mathrm{d}i\\
2\BareRoot{mp}\mathbin{.}rV\mathrm{d}t&=\mathrm{d}p.
\end{aligned}
\]
\par Nous désignerons par \(\varphi\) l’angle dont le sinus est égal à l’excentricité \(e\), de sorte que \(\BareRoot p=\BareRoot a\mathbin{.}\cos\varphi\), nous mettrons de plus pour \(m\) sa valeur \(nna^3\). \SourceLeaf{459}De cette manière nous aurons
\[
\begin{array}{rl}
61)\quad&\mathrm{d}\AscendingNode=\dfrac{arW\sin(v-\AscendingNode)}{\cos\varphi\mathbin{.}\sin i}n\mathrm{d}t\\[5pt]
62)\quad&\mathrm{d}i=\dfrac{arW\cos(v-\AscendingNode)}{\cos\varphi}n\mathrm{d}t\\[5pt]
63)\quad&\mathrm{d}p=2aar\cos\varphi\mathbin{.}Vn\mathrm{d}t.
\end{array}
\]
\par En différentiant l’équation 16, on trouve
\[
\xi\mathrm{d}x+\eta\mathrm{d}y+\zeta\mathrm{d}z=-\dfrac{m\mathrm{d}a}{2aa}\mathpunct{\cdot}
\]
\par Or les équations 4, 5, 6 donnent
\[
\begin{aligned}
(Cy-Bz)\mathrm{d}t&=xr\mathrm{d}r-rr\mathrm{d}x\\
(Az-Cx)\mathrm{d}t&=yr\mathrm{d}r-rr\mathrm{d}y\\
(Bx-Ay)\mathrm{d}t&=zr\mathrm{d}r-rr\mathrm{d}z
\end{aligned}
\]
ou bien
\[
\begin{aligned}
\mathrm{d}x&=\dfrac{x\mathrm{d}r}{r}+\dfrac{Bz-Cy}{rr}\mathrm{d}t\\[4pt]
\mathrm{d}y&=\dfrac{y\mathrm{d}r}{r}+\dfrac{Cx-Az}{rr}\mathrm{d}t\\[4pt]
\mathrm{d}z&=\dfrac{z\mathrm{d}r}{r}+\dfrac{Ay-Bx}{rr}\mathrm{d}t.
\end{aligned}
\]
\par Substituant ces valeurs il vient
\[
\begin{aligned}
-\dfrac{m\mathrm{d}a}{2aa}&=(x\xi+y\eta+z\zeta)\dfrac{\mathrm{d}r}{r}\\
&\quad+\{(Bz-Cy)\xi+(Cx-Az)\eta+(Ay-Bx)\zeta\}\dfrac{\mathrm{d}t}{rr}\\[4pt]
&=-mT\mathrm{d}r-\dfrac{m^{\frac32}\BareRoot p}{r}V\mathrm{d}t\\[4pt]
&=-\dfrac{m^{\frac32}e\sin(v-\tilde\omega)}{\BareRoot p}T\mathrm{d}t-\dfrac{m^{\frac32}\BareRoot p}{r}V\mathrm{d}t.
\end{aligned}
\]
\par Nous aurons donc
\[
\begin{array}{rl}
64)\quad&\begin{aligned}[t]\mathrm{d}a&=2a^3\operatorname{tang}\varphi\mathbin{.}\sin(v-\tilde\omega)\mathbin{.}Tn\mathrm{d}t\\&\quad+\dfrac{2a^4\cos\varphi}{r}Vn\mathrm{d}t.\end{aligned}
\end{array}
\]
Puisque \(n=\BareRoot m\mathbin{.}a^{-\frac32}\), on a \(\mathrm{d}n=-\tfrac32\dfrac{n}{a}\mathrm{d}a\), ou bien
\[
\begin{array}{rl}
65)\quad&\begin{aligned}[t]\mathrm{d}n&=-3naa\operatorname{tang}\varphi\mathbin{.}\sin(v-\tilde\omega)\mathbin{.}Tn\mathrm{d}t\\&\quad-\dfrac{3na^3\cos\varphi}{r}Vn\mathrm{d}t.\end{aligned}
\end{array}
\]
\par De plus l’équation \(p=a\cos\varphi^2\) donne \(\mathrm{d}p=\cos\varphi^2\mathbin{.}\mathrm{d}a-2a\cos\varphi\mathbin{.}\sin\varphi\mathbin{.}\mathrm{d}\varphi\); donc, tirant des équations 63, 64 les valeurs de \(\mathrm{d}p,\mathrm{d}a\), on trouve
\SourceLeaf{460}
\[
\begin{aligned}
\mathrm{d}\varphi&=aa\sin(v-\tilde\omega)\mathbin{.}Tn\mathrm{d}t\\
&\quad+\left(\dfrac{a^3\cos\varphi^2}{er}-\dfrac{ar}{e}\right)Vn\mathrm{d}t,
\end{aligned}
\]
ou à cause de \(\dfrac{a\cos\varphi^2}{r}-\dfrac{r}{a}=e(\cos(v-\tilde\omega)+\cos E)\):
\[
\begin{array}{rl}
66)\quad&\begin{aligned}[t]\mathrm{d}\varphi&=aa\sin(v-\tilde\omega)\mathbin{.}Tn\mathrm{d}t\\&\quad+aa(\cos(v-\tilde\omega)+\cos E)Vn\mathrm{d}t.\end{aligned}
\end{array}
\]
